% !TeX root = ../papers/nc-p3.tex
\section{三次接錐のPicard群と六次残余の非正規曲線}
\label{sec:residual-sextic-cubic-tangents}
\begin{theorem}\label{thm:residual-sextic-cubic-tangents}
前節の六次残余の平面外の $n\ge13$ 個の印を $P$ とする。
各印の二次Veronese評価列の間の一次関係のうち、少なくとも一つで係数が非零となる印を $P_*$ と置く。
$|P\setminus P_*|\le3$ であり、$P_*$ の全ての点は六次残余の非正規曲線上にある。
その非正規曲線の被約総次数は四以上である。
$P_*$ の三次接錐の連結Picard群は全て乗法群、または全て加法群である。
特に通常三重点（滑らかな三次接錐を持つ点）は $P$ のうち高々三個である。
\end{theorem}
\begin{proof}[二次評価の直和分解]
二次式の $P$ 上の評価は単射である。以下その双対の列の階数を用いる。
五点以下は二次式に独立な条件を課す。実際一つの印を残し他の高々四点を二組に分け、
残す点と共線な二点を同じ組に入れず、その点を避ける二平面の積を取ればよい。
各方向に他点が二つまでなのでこの分け方がある。
さらに七点以上の評価階数は六以上である。平面内で階数五以下なら非零二次式に全点が乗り、
有限窓の二次曲線上界六に反する。非平面で四点以上を持つ平面があれば、その平面内の
高々五点に対する独立性と、平面の式を掛けた一次式による外の点の分離で六に達する
（面内が四点なら外に少なくとも三点、面内が五点以上なら外に少なくとも一点）。
そのような平面がなければ任意の非共線な三点の平面と、その外の三非共線点を使って同じ結論を得る。

列の関係空間を $R\subset k^P$ とし、$\dim R=n-10\ge3$ である。互いに素な支持に関する
$R$ の直和分解があれば、対応する列の包も直和である（部分ごとの独立な列基底の和を考える）。
関係を持つ各群は六点以上で階数五以上。二群がともに関係を持てば全階数十を使い切り、他の列はなく、
いずれかが七点以上となって矛盾する。従って非零の関係は支持分割で二群に分けられない。
関係に現れない列を $c$ 本除くと、残りの階数 $r=10-c$、点数は $r+n-10\ge r+3$。
$r\le5$ なら上の評価に反する。$r=6$ なら九点以上で平面には入らない。
そのHilbert関数は次数一から $4,6,|P_*|$ である（三次評価は全射）。
非零化一次式によるArtin還元では次数二の次元が二であり、三次の次元は高々二である。
単項式初期イデアルで、二つまでの標準二次単項式だけを二次約数とする標準三次単項式は高々二つだからである。
よって $|P_*|\le8$ となり矛盾する。これで $r\ge7$, $c\le3$ を得る。
\end{proof}
\begin{proof}[三次接錐への制限]
$V$ を $P$ での三空間の吹上げ、$S$ を六次の真変換、$E$ を例外因子の和、$H_0$ を超平面の引戻しとする。
$S\sim6H_0-3E$ なので $\omega_S=(2H_0-E)|_S$。
$C=E|_S$ は相異なる平面の三次接錐 $C_i$ の非交和で、随伴により $\omega_C=\mathcal O_C(2H_0)$。
$S$ から元の六次への構造層の押出しは構造層である。実際有限点の外で同型で、
超曲面は $S_2$ であるためその外の正則関数が延びる。
従って $H^0(\mathcal O_S(2H_0))$ は入力二次式全体であり、その $C_i$ への制限は印での定数値である。
Cartier切断の双対完全列
\[
 0\longrightarrow\omega_S\longrightarrow\omega_S(C)
   \longrightarrow\omega_C\longrightarrow0
\]
から $H^0(\omega_S)=0$、また $H^0(\omega_C)\to H^1(\omega_S)$ の核は二次評価空間となる。
Serre双対性により制限 $H^1(\mathcal O_S)\to\bigoplus H^1(\mathcal O_{C_i})$ の像が
各座標の規格化を除きちょうど $R$ となる。

これは連結Picard群の制限の接写像である。標数零なので像は次元 $n-10$ の連結代数部分群である。
各平面三次の $h^0(\mathcal O)=h^1(\mathcal O)=1$ で、連結Picard群は滑らかな場合の楕円曲線、
または一次元の乗法群・加法群である。実際特異三次の被約成分の正規化は全て有理で、
正規化の単元商と冪零部分による核は線型代数群である。
これら三型の積の連結部分群の接空間は型別の直和である。連結群の最大連結アフィン正規部分は
楕円曲線積への像が零で、そのアフィン像で割った残りのアフィン群への射はアーベル多様体からの射になるため零。
可換線型群のトーラス部分とユニポテント部分も標数零で分裂する。
前段より $P_*$ はただ一型に入る。

その型が楕円なら、滑らかな例外三次の近傍では $S$ 自身が滑らかであり、
解消 $Y$ のPicard群を経由してこの次元三以上の像が得られる。よって $q(Y)\ge3$。
しかし標準跡 $f_*\omega_Y\to\omega_S$ は一般点で同型で単射なので $p_g(Y)=0$ である。
$\chi(\mathcal O_Y)<0$ から $Y$ は種数 $q(Y)$ の曲線上の線織曲面に双有理である。
楕円例外曲線からこの底への射は定数だから、底から引き戻される $\operatorname{Pic}^0(Y)$ の制限は全て零である。
この矛盾で楕円型を除いた。
\end{proof}
\begin{proof}[非正規軌跡への支持]
連結Picard群のアフィン部分は上の像全体に全射する。実際その像で割るとアーベル多様体から
アフィン群への準同型が残るだけである。$S$ の正規化を $\nu:T\to S$ とする。
正規射影多様体の $\operatorname{Pic}^0(T)$ はアーベル多様体なのでこのアフィン部分は $T$ で自明になる。
引戻しの核は単元層の完全列により
$H^0(S,\nu_*\mathcal O_T^\times/\mathcal O_S^\times)$ である（両射影整多様体の大域単元は定数）。
この群は非正規軌跡の互いに素な閉成分群ごとに直積に分かれ、それぞれの境界準同型による線束は
その支持の外で自明である。この分解は基礎環を二重数としても同じで、接空間でも成り立つ。

もし $p\in P_*$ で元の六次が正規なら、その近傍で $S$ が非正規になる箇所は $C_p$ 上だけであり、
この閉部分は他の非正規軌跡から分離される。$C_p$ で非零の接方向を与える上の単元商の元を
この部分だけに制限し、他を恒等とすれば、$R$ に $p$ 一点だけを支持とする非零元を得てしまう。
二次評価の列は零でないので不可能である。こうして $P_*$ は元の非正規軌跡に入る。

超曲面の非正規台は二重以上の曲線の和である。これが次数三以下なら、平面的な曲線には七印、
複数の直線・二次曲線には合計九印までしか載らない。十印以上の $P_*$ を含むには非平面整三次、
すなわち有理正規三次 $J$ でなければならない。
六次式の $J$ に沿う二次の初期項は
$\operatorname{Sym}^2(\mathcal I_J/\mathcal I_J^2)\otimes\mathcal O_J(6)$ の非零切断である。
零なら三重曲線が生じて前定理に反する。有理正規三次の法束は $\mathcal O_{\mathbb P^1}(5)^2$ だから
初期項は三本の次数八の切断になる。この法束の計算は、三つの標準二次生成元を $J$ に制限した
余法表示でも得られる。二つの線型Hankel関係の像は $J$ 上で原始な二次列の張る直線となり、
$\mathcal O(-8)\to\mathcal O(-6)^3$ の余核の双対は、二項全二次系の二つの一次Koszul型関係により
$\mathcal O(5)^2$ である。各三重印では二次初期項は零となるので、十印は収容できない。
\end{proof}
